\section{Termination}\label{sec:termination}

In the previous section we show how tableaux yield a satisfiability calculus.
In this section, we prove that we can restrict these tableaux to be finite, without loosing completeness.
Such a result establishes that the satisfiability problem for $\Hxpath$ is decidable.

We begin by observing that the analyticity of the tableau ensures that all expressions introduced during expansion are constructed from subcomponents of the root, reducing the search space. 
The main technical challenge lies in managing the introduction of new nominals and new accessibility constraints through applications of the ($\Diamond$) and (child) rules. 
To resolve this, we adapt the termination argument for Hybrid Logic tableaux in~\cite{BolanderB07}. In particular, we impose a controlled strategy for applying the ($\Diamond$) and (child) rules to ensure that the resulting structure remains tree-like. 

\begin{defi}
  \label{def:atform}
  Let $\Theta$ be a branch of a tableau with root $i{:}\varphi$.
  For $j \in \noms(\Theta)$, define:
  \begin{align*}
    \at_{\nexpr}^{\Theta}(j) & = \setof{\psi}{{j{:}\psi} \in (\Theta \cap \qsub(i{:}\varphi))} \\
    \at_{\pexpr}^{\Theta}(j) & = \setof{\alpha}{\pm \<@_j\alpha \cmpr \beta \> \in (\Theta \cap \qsub(i{:}\varphi))}.
  \end{align*}
  where $\pm$ indicates that $\<@_j\alpha \cmpr \beta \>$ may or may not be preceded by a negation.
  In addition, define $\at^\Theta(j) = \at_{\nexpr}^{\Theta}(j) \cup \at_{\pexpr}^{\Theta}(j)$.
\end{defi}

Intuitively, the set $\at_{\nexpr}^{\Theta}(j)$ captures the subcomponents of the root of the tableau true at $j$ in the branch.
In turn, the set $\at_{\pexpr}^{\Theta}(j)$ captures the subpaths present in the data (in)equalities in the root of the tableau that start at $j$ in the branch.
%We drop the superscript $\Theta$ when it can be understood from the context.

%\bigval{include an example for $\at_{\pexpr}^{\Theta}(j)$}
\begin{exa} 
  \label{ex:at}
  Let $\Theta$ be the sole branch in the tableau in \Cref{ex:usetableaux}.
  Then
%    \begin{align*}
%      \at^{\Theta}(5) = {} &
%      {\at_{\nexpr}^{\Theta}(5) \cup \at_{\pexpr}^{\Theta}(5)} \\&
%      {\set{\tup{@_2\astep[b]2? =_\cmp \astep[b](q \land 3)?}} \cup \set{@_2\astep[b]2?,\astep[b](q\land 3)?}} \\&
%      \set{\tup{ @_2\astep[b]2? =_\cmp \astep[b](q \land 3)?},@_2\astep[b]2?,\astep[b](q\land 3)?}.
%    \end{align*}
$$	      \at^{\Theta}(5) =       {\at_{\nexpr}^{\Theta}(5) \cup \at_{\pexpr}^{\Theta}(5)} =	      {\set{\tup{@_2\astep[b]2? =_\cmp \astep[b](q \land 3)?}} \cup \set{@_2\astep[b]2?,\astep[b](q\land 3)?}}.$$
  Notice how the set $\at^{\Theta}(5)$ gives us a hint of the rules which can be applied to the node expressions having $5$ as a prefix. 
  We retake this idea later on.
\end{exa}

The set $\qsub(i{:}\varphi)$ of quasi subcomponents of the label of the root of a tableau is infinite, since its considers all possible ways of prefixing it with nominals.
However, we can establish that in fact $\at^\Theta$ only considers those  quasicomponents that are made of either a subcomponent or a negated subcomponent, i.e., it is restricted to nominals appearing in a branch of the tableau. This is stated in the following result.

\begin{prop} \label{prop:at:finite}
  For every branch $\Theta$, the set $\at^\Theta(j)$ is finite.
\end{prop}
\begin{proof}
  Let $\Theta$ be a branch of a tableau with root $i{:} \varphi$.
  From \Cref{prop:analytic}, we get: 
  \begin{align*}
    {\at^{\Theta}_{\nexpr}}(j)
      & \subseteq  \setof{\psi, \lnot \psi}{\psi \in \sub(i{:}\varphi) } 
    \\
    \at^{\Theta}_{\pexpr}(j)
      & \subseteq \setof{\alpha}{{\pm}\tup{\alpha \cmpr \beta} \in \sub(i{:}\varphi)}.
  \end{align*}
  Since the set $\sub(i{:}\varphi)$ is finite, we immediately get that the set $\at^\Theta(j)$ is finite.
\end{proof}

\Cref{prop:at:finite} tells us that the set of expressions that are true at a given node must be finite. Still, there is no guarantee that every tableau must be finite. It could be the case that the tableau is infinitely branching, or that it has a branch of infinite size. We prove that neither is the case for our tableaux.

\begin{defi} \label{def:generated}
  Let $\Theta$ be a branch of a tableau, and $\{i, j\} \subseteq \noms(\Theta)$.
  We say that \emph{$j$ is generated by~$i$}, and write $i \prec_\Theta j$, iff there is $\pathh\in\rels$ such that $i{:}\tup{\pathh}j\in\Theta$ is an accessibility constraint. 
  Finally, we say that $i$ is an \emph{ancestor} of $j$, 
    and write $i \prec_\Theta^{+} j$,
  iff $(i,j)$ is in the transitive closure of $\prec_\Theta$.
\end{defi}

The following lemma tells us that every nominal in the branch of a tableau generates a finite number of new nominals. More precisely, this result tells us that applications of ($\Diamond$) and (child) result in a finite set of trees, each of which has a finite number of branches.

\begin{lem} \label{lemma:graph}
  Let $\Theta$ be a branch of a tableau.
  The graph $G_{\Theta} = (\noms(\Theta), \prec_\Theta)$ is a finite set of finitely branching trees.
\end{lem}
\begin{proof}
  That $G_{\Theta}$ is a finite set of trees follows from the following three facts.
  First, for every $\{i,j,k\} \subseteq \noms(\Theta)$, if $j \prec_\Theta i$ and $k \prec_\Theta i$, then, $j = k$, i.e., if a nominal is generated by another nominal, then it is generated by at most one other nominal.
  Second, the set of ancestors of a nominal $i \in \noms(\Theta)$ that is generated by another nominal has a least element.
  Third, for any $i \in \noms(\Theta)$, the set $\at^\Theta(i)$ is finite (see~\Cref{prop:at:finite}), this implies that there is at most $\card(\at^\Theta(i))$ new nominals generated from $i$.
\end{proof}

Now that we know that applications of ($\Diamond$) and (child) always result in a finite set of finitely branching trees, we can move on to show that every branch in such trees is of finite size. Such a result establishes that tableaux are finite. 

\begin{lem} \label{lemma:chain{:}infinite}
  A branch of a tableau is infinite iff there is an infinite chain $i \prec_\Theta j  \prec_\Theta \dots$.
\end{lem}

\begin{proof}
  The right-to-left direction is trivial.
  To prove the left-to-right direction, let  $\Theta$ be an infinite branch.
  From \Cref{prop:at:finite}, the branch $\Theta$ must contain infinitely many distinct nominals.
  This implies that $\noms(\Theta)$ in \Cref{lemma:graph} must be infinite.
  Since, from \Cref{lemma:graph}, we know that $G_\Theta$ is a finite set of finitely branching trees, it must be that some tree in $G_\Theta$ contains an infinite chain $i \prec_\Theta j \prec_\Theta \dots $.
\end{proof}

\Cref{def:maxsize} introduces the last piece needed to prove termination.

\begin{defi}\label{def:maxsize}
  Let $\Theta$ be a branch of a tableau, and let $i \in \noms(\Theta)$. Define
  \begin{align*}
    \maxnode(i) & = 
      \max\setof{ \size(\varphi)}{ \varphi \in \at^\Theta_{\nexpr}(i) } \\
    \maxpath(i) &=
      \begin{cases}
        \max\setof{\size(\alpha)}{\alpha \in \at^{\Theta}_{\pexpr}(i)} & \text{if } \at^{\Theta}_{\pexpr}(i) \neq \emptyset \\
        0 & \text{otherwise}\\
      \end{cases}\\
    \maxsize(i)
& = \max (\maxnode(i) , \maxpath(i)).
  \end{align*}
  % where
  % \begin{align*}
  %   \maxnode(i) & = 
  %     \max\setof{ \size(\varphi)}{ \varphi \in \at^\Theta_{\nexpr}(i) }
  %   &
  %   \maxpath(i) &=
  %     \begin{cases}
  %       \max\setof{\size(\alpha)}{\alpha \in \at^{\Theta}_{\pexpr}(i)} & \text{if } \at^{\Theta}_{\pexpr}(i) \neq \emptyset \\
  %       0 & \text{otherwise}.
  %     \end{cases}
  % \end{align*}
\end{defi}

Intuitively, $ \maxsize(i)$ indicates the size of the largest node expression true at $i$ in $\Theta$, or the length of the longest path expression starting at $i$.
The idea is to prove that such a measure of size decreases after the application of the tableaux rules. 

\begin{lem} \label{lemma:decreasing:length}
  Let $\Theta$ be a saturated branch. If $i \prec_\Theta j$, then $\maxsize(i) > \maxsize(j)$.
\end{lem}
\begin{proof}
 % Recall the definition of $\size$ in \Cref{def:size}.
  Let $i \prec_\Theta j$.
  We have three possible cases:
    (1) $\maxsize(j) = 0$,
    (2) $\maxsize(j) = \size(\psi)$ for some node expression $\psi$, and
    (3) $\maxsize(j) = \size(\alpha)$ for some path expression $\alpha$.
  The latter has two possible sub-cases:
    (a) $\tup{@_j\alpha \cmpr \beta} \in \Theta$ and $\lnot\tup{@_j\alpha \cmpr \beta} \notin \Theta$, and
    (b) $\lnot\tup{@_j\alpha \cmpr \beta} \in \Theta$ and $\tup{@_j\alpha \cmpr \beta} \notin \Theta$.
  The result is immediate for (1). 
  We prove the remaining cases below.
  
  \begin{description}
    \item[(Case 2)]
      Let $\maxsize(j) = \size(\psi)$ for some node expression $\psi$.
      We need to prove $\maxsize(i) > \size(\psi)$.
      % It is clear that $j{:}\psi$ cannot result from ($\lnot \lnot$), ($\land$), ($\lnot \land$), (nom), ($\lnot$nom), or (copy).
      % The first three rules imply that $\size(\psi) > \maxsize(j)$ and this is a contradiction.
      % The last three rules cannot be applied since their side conditions are not met.
      % The remaining rules are:
      %
      % \begin{description}
        % \item[($j{:}\psi$ results from ($\lnot \lnot$), ($\land$), or ($\lnot \land$))]
        % We get a contradiction from $\maxsize(j) > \size(\psi)$.
        % This means $j{:}\psi$ cannot be obtained from these rules.
        %
        % \item[($j{:}\psi$ results from (nom) or ($\lnot$nom))]
        % We get a contradiction since $i \prec_\Theta j$ and thus $j$ must be new.
        % This means $j{:}\psi$ cannot be obtained from these rules.
      It is clear that $j{:}\psi$ can only result from ($\Diamond$), ($\lnot\Diamond$), (ref), (sym), or (trans).
        % \item[($\Diamond$) and ($\lnot\Diamond$)]
        If $j{:}\psi$ results from ($\Diamond$), using \Cref{lemma:graph}, we know $i{:}\tup{\child}\psi \in \Theta$.
        This implies, $\maxsize(i) \geq \size(\tup{\child} \psi) > \size (\psi) = \maxsize(j)$ as needed.
        %
        In turn, if $j{:}\psi$ results from ($\lnot\Diamond$), we know $\psi = \lnot\psi'$.
        Using \Cref{lemma:graph} again, we know $\{i{:}\lnot\tup{\child}\psi', i{:}\tup{\pathh}j\} \subseteq \Theta$.
        This implies, $\maxsize(i) \geq \size(\lnot\tup{\child} \psi') > \size (\psi) = \maxsize(j)$ as needed.
        %
        % \item[(ref), (sym), and (trans)]
        Lastly, if $j{:}\psi$ results from (ref), (sym), or (trans), we have ${\size(\psi)=1}$. 
        At the same time, from $i \prec_{\Theta}j$, we have $\{i{:}\<\child\>\chi, \<\goto{i}\child\alpha \cmpr \beta\>\} \cap \Theta \neq \emptyset$.
        Immediately, $\maxsize(i) > 1 = \size(\psi)$ as needed.

        % \item[($j{:}\psi$ results from (copy))]
        % We have $j < k$.
        % This implies $k$ is not a root nominal (see~\Cref{def:tableau}).
        % But this is a contradiction; since in an application of (copy) $k$ must be a root nominal.
      % \end{description}
  
    \item[(Case 3.a)]
      Let $\maxsize(j) = \size(\alpha)$, $\tup{@_j\alpha \cmpr \beta} \in \Theta$, and $\lnot\tup{@_j\alpha \cmpr \beta} \notin \Theta$.
      We need to prove $\maxsize(i) > \size(\alpha)$.
      % It is clear that $\tup{@_j\alpha \cmpr \beta}$ cannot result from (int$_1$), (${\lnot}{\cmpr}$), (dRef), (test), ($@$), ($\cup$), (com$_1$), or (dTrans).
      % % All these cases break the maximality of $\maxsize(j)$ or fail to comply with the applicability conditions of the rule.
      % The remaining rules are:
      %
      % \begin{description}
      %   \item[(child)]
      It is clear that $\tup{@_j\alpha \cmpr \beta}$ can only result from (child).
      Using \Cref{lemma:graph}, we know that $\tup{@_i\pathh\alpha \cmpr \beta} \in \Theta$.
      Then, $\maxsize(i) \geq \size(\tup{@_i\pathh\alpha \cmpr \beta}) > \size (\tup{\alpha \cmpr \beta}) = \maxsize(j)$ as needed.
    
    \item[(Case 3.b)]
      Let $\maxsize(j) = \size(\alpha)$, $\lnot\tup{@_j\alpha \cmpr \beta} \in \Theta$, and $\tup{@_j\alpha \cmpr \beta} \notin \Theta$.
      We need to prove $\maxsize(i) > \size(\alpha)$.
      % It is clear that $\tup{@_j\alpha \cmpr \beta}$ cannot result from (int$_1$), (${\lnot}{\cmpr}$), (dRef), (test), ($@$), ($\cup$), (com$_1$), or (dTrans).
      % % All these cases break the maximality of $\maxsize(j)$ or fail to comply with the applicability conditions of the rule.
      % The remaining rules are:
      %
      % \begin{description}
      %   \item[(child)]
      It is clear that $\lnot\tup{@_j\alpha \cmpr \beta}$ can only result from ($\lnot$child).
      Using \Cref{lemma:graph}, we know that $\lnot\tup{@_i\pathh\alpha \cmpr \beta} \in \Theta$.
      Then, $\maxsize(i) \geq \size({\lnot\tup{@_i\pathh\alpha \cmpr \beta}}) > \size (\lnot\tup{\alpha \cmpr \beta}) = \maxsize(j)$ as needed. \qedhere

    % \item[(Case 2.b)]
    % We need to prove $\maxsize(i) > \size(\alpha)$.
    % It is clear that $\tup{@_j\alpha \cmpr \beta}$ cannot result from (int$_1$), (${\lnot}{\cmpr}$), (dRef), (test), ($@$), or ($\cup$).
    % All these cases break the maximality of $\maxsize(j)$ or fail to comply with the applicability conditions of the rule.
    % The remaining cases are:

    % \begin{description}
    %   \item[(child)]
    %   In the case of (child), we know that $\tup{@_k\pathh\alpha \cmpr \beta} \in \Theta$ and $k \prec_\Theta j$.
    %   At the same time, we have assumed $i \prec_\Theta j$.
    %   From \Cref{lemma:graph}, every non-root nominal is generated by at most one other nominal.
    %   This implies $k = i$.
    %   Then, $\maxsize(i) \geq \size(\tup{\pathh\alpha \cmpr \beta}) > \size (\tup{\alpha \cmpr \beta}) = \maxsize(j)$.
    %   In the case of ($\lnot$child), we know $\{\lnot\tup{@_i\pathh\alpha \cmpr \beta}, i{:}\tup{\child}j\} \subseteq \Theta$.
    %   Then, $\maxsize(i) \geq \size(\lnot\tup{\child} \psi') > \size (\psi) = \maxsize(j)$.

    %   \item[($\tup{@_j\alpha \cmpr \beta}$ results from (com$_1$) or (dTrans))]
    %   Similar to the case (sym) and (trans), respectively.
    % \end{description}
  \end{description}
\end{proof}

\begin{lem} \label{prop:termination}	
  Let $\varphi$ be a node expression.
  Any tableau for $\varphi$ is finite.
\end{lem}
\begin{proof}
  The proof is by contradiction. Suppose there is an infinite tableau for $\varphi$. Then it must be that this tableau has an infinite branch $\Theta$.
  By \Cref{lemma:chain{:}infinite}, there exists an infinite chain
  $
   i \prec_\Theta j \prec_\Theta \dots  
  $
  At the same time, by \Cref{lemma:decreasing:length}, we have 
  $
    \maxsize(i) > \maxsize(j) > \dots  
  $
  This yields a contradiction; since for any nominal $i$, it its trivial that $\maxsize(i) \geq 0$.
\end{proof}

Combining the result of \Cref{prop:termination} with the completeness of our tableau-based satisfiability procedure, we obtain the main result of this section.

\begin{thm}\label{th:decidability}
  The satisfiability problem for $\Hxpath$ is decidable.
\end{thm}
% \begin{proof}
%   Let $\varphi$ be a node expression. We can list the tableaux for $\varphi$ as if they where build in a depth-first search manner.
%   For every position in the list, we can also explore the tableau in that position and check if it is closed, or if it contains an open and saturated branch.
%   Since we know from \Cref{prop:termination} that tableaux for $\varphi$ are finite, both listing and exploring have intuively defined algorithms.
%   This establishes that the satisfiability problem for $\Hxpath$ is decidable.
% \end{proof}

In \Cref{sec:complexity} we refine the result in \Cref{th:decidability} and characterize the exact complexity of the satisfiability problem for $\Hxpath$. 